\documentclass[10pt,a4paper]{amsart}
\usepackage[margin=19mm]{geometry}
\usepackage{amsmath,amssymb,mathtools}
\usepackage[scr=rsfs,cal=euler]{mathalfa}
\usepackage{xfrac}
\usepackage[unicode,hidelinks,bookmarks=false]{hyperref}
\newcommand{\ie}{i.e.\ }
\newcommand{\cf}{cf.\ }
\newcommand{\resp}{respectively\ }
\newcommand{\Subsection}{\subsection}
\providecommand{\nobreakdash}{\nobreak\hbox{-}}
\setlength{\emergencystretch}{2em}
\multlinegap0pt
\usepackage{amssymb}
\usepackage{mathtools}
\usepackage{tikz}
\newtheorem{proposition}{Proposition}
\title{\textbf{On $\mathscr{M}$-arrangements of conics and lines with ordinary singularities}}
\author{Marek Janasz, Piotr Pokora}
\date{Source: arXiv:2512.24707v2 (CC BY 4.0)}
\begin{document}
\maketitle

\noindent\emph{Source and licence.} \textbf{On $\mathscr{M}$-arrangements of conics and lines with ordinary singularities}. Version arXiv:2512.24707v2 (CC BY 4.0), distributed by arXiv under the Creative Commons Attribution 4.0 International licence, http://creativecommons.org/licenses/by/4.0/, as recorded in the arXiv licence metadata for this exact version and shown on the abstract page https://arxiv.org/abs/2512.24707v2. Full licence notice: \texttt{licenses/arxiv-2512.24707-CC-BY-4.0.txt}.\par\smallskip

\noindent\textit{Editorial scope.} Counted unit: the article's proposition proposition (source0.tex lines 676--685). The article's complete original proof of it is delivered with it (source0.tex lines 686--695, one \texttt{proof} environment of 10 lines). No mathematics is written, repaired or re-derived: the statement and the proof are the article's own lines, in the article's own notation, and no retained line is substituted or re-flowed.  No further source-local context is needed: every cross-reference of every retained line resolves inside the two delivered spans.  Nothing was omitted from any retained span, so the package carries no omission record.  No retained line contains a citation, so the wrapper carries no bibliography and no bibliography entry.  No retained line contains a figure environment, an image-inclusion command or drawing code, so no image is delivered and the package declares no asset dependency.  Editorial formatting only, in addition: an AMS \texttt{amsart} preamble that restates the article's own declarations and macros used by the retained lines, namely the environments \texttt{proposition} on separate counters, together with the standard packages those lines need.  The retained environments print in this document as Proposition 1 in order of appearance, whereas the article prints the same items with the numbering of its own full document.  The complete native source of the article as distributed in the arXiv e-print for this exact version is bundled unchanged as \texttt{sources/191922/source0.tex}.
\begin{proposition}
Let $C = \{f=0\}\subset \mathbb{P}^{2}_{\mathbb{C}}$ be an $\mathscr{M}$-curve of degree $d$. Then
$${\rm reg}(M(f)) = \begin{cases}
3m-2 & \text{ if } d=2m\geq 6, \text{ or}\\
3m-1 & \text{ if } d=2m+1 \geq 5
\end{cases}$$
and
$${\rm reg}({\rm AR}(f)) = m+1$$
both for $d=2m\geq 6$ or $d=2m+1\geq 5$.
\end{proposition}

\begin{proof}
It is enough to observe that 
$${\rm reg}({\rm AR}(f)) = {\rm st}(f) - d + 3 = d_{2},$$
and the claim follows from the fact that, if $C = \{f=0\} \subset \mathbb{P}^{2}_{\mathbb{C}}$ is an $\mathscr{M}$-curve of degree $d$, then $C = \{f=0\}$ is free with exponents $(d_{1},d_{2})$, where
\begin{align*}
(d_{1},d_{2}) = (m-2, m+1) & \text{ if } d=2m\geq 6, \text{ or} \\
(d_{1},d_{2}) = (m-1, m+1) & \text{ if } d=2m+1\geq 5.
\end{align*}
This gives us ${\rm reg}({\rm AR}(f)) = m+1$, and then the values for ${\rm reg}(M(f))$ follow.
\end{proof}

\end{document}
